% =============================================================================
% SECTION 4 — EXPERIMENTS
% =============================================================================
% Notes:
%  - The main table uses a multicolumn header row per base model (rather than
%    multirow), which avoids the \multirow + \rowcolor interaction that hid
%    the 7B label in the previous version.
%  - The light-blue Δ row requires \usepackage[table]{xcolor} (or colortbl).
%  - All numerical results are means over three independent seeds; std bands
%    in red are placeholders awaiting final replication runs.
% =============================================================================

\section{Experiments}
\label{sec:exp}

\subsection{Setup}
\label{sec:exp:setup}

\paragraph{Benchmarks.}
We evaluate on three groups of benchmarks. \emph{Coding agent benchmarks}
comprise SWE-Bench-Verified (500 instances) and SWE-Bench-Lite
(300 instances)~\citep{swebench}; both are in-domain for the
post-training pipelines and constitute our primary measure of end-task
progress. \emph{Non-agent coding benchmarks} cover
LiveCodeBench~\citep{livecodebench}, OJBench~\citep{ojbench}, and
FullStackBench-EN~\citep{fullstackbench}; these probe pure code generation
without an agent loop and serve as a regression check for capability
erosion. \emph{Tool-use and out-of-domain agent benchmarks} include
Terminal-Bench~\citep{terminalbench}, $\tau$-bench~\citep{taubench}, and
BFCL~\citep{bfcl}; the latter two contain no Python code-editing
trajectories at all, and they are the appropriate stage on which to test
our central hypothesis that the function-call inductive bias transfers
across task families.

\paragraph{Mid-training and post-training pipeline.}
We mid-train Qwen2.5-Coder-Instruct (7B, 14B, 32B) and Qwen3-8B on the
selected FIM corpus. The training objective is the standard FIM loss on
the middle span (rationale plus body), packed to the model's native
context length and using the model's native FIM sentinel tokens to
preserve compatibility with any FIM signal already learned during
pretraining. After mid-training we apply existing agentic post-training
pipelines without modification: R2E-Gym~\citep{r2egym} or
SWE-Smith~\citep{swesmith} for the Qwen2.5-Coder runs, and
SWE-Lego~\citep{swelego} for Qwen3-8B. \emph{Important deviation for
Qwen3-8B}: we train for $2$ epochs rather than the $4$ epochs used in the
original SWE-Lego setting, as the latter exhibited noticeable overfitting
in our setup. Hyperparameters, token budgets, and learning-rate schedules
are reported in Appendix~\ref{app:hp}.

\paragraph{Evaluation protocol.}
\textbf{All numerical results are reported as the mean over three
independent evaluation seeds}; the ``\%'' symbol is omitted inside table
cells. Standard-deviation bands shown in red (\textcolor{red}{$\pm$?}) are
placeholders awaiting final replication runs. Results are obtained on the
\emph{final} checkpoint of the training pipeline. A ``baseline'' row
refers to the base model followed by an agentic post-training pipeline,
while ``ours'' refers to the same base followed by FIM mid-training and
then by the identical post-training pipeline. We do not directly evaluate
mid-training-only checkpoints, because FIM data without subsequent
instruction-style post-training degrades instruction-following ability and
makes any comparison against instruction-tuned baselines unfair on
instruction-following benchmarks. Consequently, every gain we attribute
to mid-training should be read as a gain that survives subsequent
post-training. The agent harness is fixed by the post-training pipeline:
R2E-Gym runs use the R2E-Gym scaffold (an OpenHands
fork)~\citep{openhands, r2egym}; SWE-Smith runs use
SWE-agent~\citep{sweagent}; the Qwen3-8B runs use OpenHands directly,
because the SWE-Lego post-training data exceeds the 32K context length of
the R2E-Gym/SWE-Smith scaffolds while remaining within Qwen3's
40{,}960-token context. Maximum turn limits, decoding temperature, and
per-step timeouts are deferred to Appendix~\ref{app:hp}.

\paragraph{Models and baselines.}
For Qwen2.5-Coder-Instruct at 7B / 14B / 32B~\citep{qwencoder} we report
(i) the instruction-tuned base without any agentic training, (ii)
post-training with R2E-Gym reproduced under our setup, (iii) the official
R2E-Gym numbers (in grey, where available), (iv) our recipe of FIM
mid-training followed by R2E-Gym post-training, and at 7B additionally
the analogous SWE-Smith~\citep{swesmith} variants. For Qwen3-8B we follow
the SWE-Lego~\citep{swelego} pipeline; the cross-base-model comparison
thus also varies the post-training pipeline, a confound we discuss in
Section~\ref{sec:exp:main}.

% =============================================================================
% 4.2 Main Results
% =============================================================================
\subsection{Main Results on Agent Benchmarks}
\label{sec:exp:main}

% -----------------------------------------------------------------------------
% Main table: uses multicolumn headers per base model (avoids the
% multirow+rowcolor interaction that previously hid the 7B label).
% Δ rows are highlighted via \rowcolor{blue!8} (requires
% \usepackage[table]{xcolor}).
% -----------------------------------------------------------------------------
\begin{table}[t]
\centering
\caption{Main results on coding agent benchmarks (SWE-Bench-Verified and
SWE-Bench-Lite). All cells are percentages, reported as the mean over
three independent seeds. Bolded rows are our recipe (FIM mid-training
followed by the same post-training pipeline as the row above). Greyed
rows are reproduced from officially released numbers; black rows are
reproduced under our setup. \textcolor{blue!70!black}{Light-blue rows
report the absolute gain $\Delta$ of our recipe over the same
post-training pipeline reproduced under our setup.} Std bands shown in
red are placeholders awaiting final replication results.}
\label{tab:main}
\small
\setlength{\tabcolsep}{6pt}
\begin{tabular}{lcc}
\toprule
Setting & SWE-Bench-Verified & SWE-Bench-Lite \\
\midrule
\multicolumn{3}{l}{\textit{Qwen2.5-Coder-7B-Instruct}} \\
\quad --- (no agentic training)              & 1.8\,($\pm$1.3)                              & 1.0\,($\pm$1.0) \\
\quad + R2E-Gym (reproduced)                 & 15.00\,(\textcolor{red}{$\pm$?})              & 11.33\,(\textcolor{red}{$\pm$?}) \\
\quad + R2E-Gym (official)                   & \textcolor{gray}{19.00}                        & \textcolor{gray}{11.00} \\
\quad \textbf{+ FIM-Midtrain + R2E-Gym}      & \textbf{17.80}\,(\textcolor{red}{$\pm$?})     & \textbf{15.00}\,(\textcolor{red}{$\pm$?}) \\
\rowcolor{blue!8}
\quad\quad $\Delta$ (vs.\ reproduced)        & $+2.80$                                        & $+3.67$ \\
\cmidrule(lr){1-3}
\quad + SWE-Smith (reproduced)               & 12.30\,(\textcolor{red}{$\pm$?})              & 14.20\,(\textcolor{red}{$\pm$?}) \\
\quad + SWE-Smith (official)                 & \textcolor{gray}{15.20}                        & \textcolor{gray}{11.70} \\
\quad \textbf{+ FIM-Midtrain + SWE-Smith}    & \textbf{17.60}\,(\textcolor{red}{$\pm$?})     & \textbf{14.70}\,(\textcolor{red}{$\pm$?}) \\
\rowcolor{blue!8}
\quad\quad $\Delta$ (vs.\ reproduced)        & $+5.30$                                        & $+0.50$ \\
\midrule
\multicolumn{3}{l}{\textit{Qwen2.5-Coder-14B-Instruct}} \\
\quad --- (no agentic training)              & 4.0\,($\pm$1.6)                              & 2.7\,($\pm$1.0) \\
\quad + R2E-Gym (reproduced)                 & 26.20\,(\textcolor{red}{$\pm$?})              & 18.00\,(\textcolor{red}{$\pm$?}) \\
\quad + R2E-Gym (official)                   & \textcolor{gray}{26.80}                        & \textcolor{gray}{20.67} \\
\quad \textbf{+ FIM-Midtrain + R2E-Gym}      & \textbf{29.20}\,(\textcolor{red}{$\pm$?})     & \textbf{22.00}\,(\textcolor{red}{$\pm$?}) \\
\rowcolor{blue!8}
\quad\quad $\Delta$ (vs.\ reproduced)        & $+3.00$                                        & $+4.00$ \\
\midrule
\multicolumn{3}{l}{\textit{Qwen2.5-Coder-32B-Instruct}} \\
\quad --- (no agentic training)              & 7.0\,($\pm$1.3)                              & 3.0\,($\pm$1.4) \\
\quad + R2E-Gym (reproduced)                 & 31.80\,(\textcolor{red}{$\pm$?})              & 24.70\,(\textcolor{red}{$\pm$?}) \\
\quad + R2E-Gym (official)                   & \textcolor{gray}{34.40}                        & \textcolor{gray}{23.77} \\
\quad \textbf{+ FIM-Midtrain + R2E-Gym}      & \textbf{35.10}\,(\textcolor{red}{$\pm$?})     & \textbf{26.80}\,(\textcolor{red}{$\pm$?}) \\
\rowcolor{blue!8}
\quad\quad $\Delta$ (vs.\ reproduced)        & $+3.30$                                        & $+2.10$ \\
\midrule
\multicolumn{3}{l}{\textit{Qwen3-8B-Base}} \\
\quad --- (no agentic training)              & 7.60\,(\textcolor{red}{$\pm$?})               & 5.80\,(\textcolor{red}{$\pm$?}) \\
\quad + SWE-Lego                             & 31.80\,(\textcolor{red}{$\pm$?})              & 27.30\,(\textcolor{red}{$\pm$?}) \\
\quad \textbf{+ FIM-Midtrain + SWE-Lego}     & \textbf{35.00}\,(\textcolor{red}{$\pm$?})     & \textbf{32.70}\,(\textcolor{red}{$\pm$?}) \\
\rowcolor{blue!8}
\quad\quad $\Delta$ (vs.\ reproduced)        & $+3.20$                                        & $+5.40$ \\
\bottomrule
\end{tabular}
\end{table}

Table~\ref{tab:main} summarizes the in-domain agent results.

\paragraph{Mid-training delivers consistent gains across model scales.}
Holding the post-training pipeline fixed at R2E-Gym, FIM mid-training
improves SWE-Bench-Verified by $+2.80$ at 7B, $+3.00$ at 14B, and $+3.30$
at 32B. The absolute gain stays within a narrow band as the base model
grows by more than $4{\times}$ in parameters. SWE-Bench-Lite shows the
same direction at every scale ($+3.67$ / $+4.00$ / $+2.10$). We do not
interpret this as evidence of monotone scaling, but rather as evidence
that the structural prior is not absorbed by larger pretrained models in
the practical deployment range we evaluate. We do not extrapolate these
trends beyond 32B.

\paragraph{The improvement transfers across post-training pipelines.}
Replacing R2E-Gym with SWE-Smith on the same 7B base, FIM mid-training
yields a gain of $+5.30$ points on SWE-Bench-Verified, larger than the
$+2.80$ observed under R2E-Gym. On SWE-Bench-Lite the SWE-Smith pairing
gains only $+0.50$, smaller than the $+3.67$ for the corresponding
R2E-Gym pairing; the cross-pipeline transfer is therefore clearer on
Verified than on Lite. Taken together, the two pipelines indicate that
mid-training is not specifically tuned to a single post-training data
distribution, but the magnitude of its benefit can depend on which
baseline pipeline it is composed with.

\paragraph{The improvement transfers across base-model families.}
Switching to Qwen3-8B and the SWE-Lego post-training pipeline,
mid-training improves SWE-Bench-Verified by $+3.20$ points and
SWE-Bench-Lite by $+5.40$ points, comparable to the gains observed on
Qwen2.5-Coder. We note that this comparison varies the post-training
pipeline simultaneously with the base model (SWE-Lego is required because
Qwen3-8B's $40{,}960$-token context exceeds that of the
R2E-Gym/SWE-Smith scaffolds), so the cross-base-model result should be
read as ``the recipe is not specific to the Qwen2.5-Coder family under
one alternative agentic pipeline,'' rather than as a guarantee that it
generalizes across all base-model families. We elaborate this caveat in
Section~\ref{sec:limitations}. The full Qwen2.5-Coder scaling curves
($7$B/$14$B/$32$B) are given in Appendix~\ref{app:scaling}.

% =============================================================================
% 4.3 Capability preservation and cross-domain transfer
% =============================================================================
\subsection{Capability Preservation and Cross-Domain Transfer}
\label{sec:exp:gen}

A natural concern with any post-training stage that specializes a model
toward a narrow agent task is that it may erode capabilities the base
model already possessed. We test for this directly on six benchmarks that
sit outside the SWE-Bench distribution and, in the case of $\tau$-bench
and BFCL, outside the coding domain entirely. Because this analysis
requires a controlled comparison among (instruct base) vs.\ (post-training
only) vs.\ (mid-training $+$ post-training), and the corresponding
controlled triple is the most expensive to produce, we conduct it at the
14B scale with the R2E-Gym pipeline. Results are summarized in
Table~\ref{tab:cap}.

% -----------------------------------------------------------------------------
% Capability preservation table
% -----------------------------------------------------------------------------
\begin{table}[t]
\centering
\caption{Capability preservation and cross-domain transfer at 14B with the
R2E-Gym post-training pipeline. ``Non-agent coding'' tests pure code
generation; ``Agent OOD'' is closer to coding agents but in a different
environment; ``Tool use'' tests non-coding tool-call behavior. Bold marks
the better trained variant per column (post-only vs.\ ours); the Instruct
row is shown as a reference ceiling and is not in competition. All cells
are percentages, std bands in red are placeholders.}
\label{tab:cap}
\small
\setlength{\tabcolsep}{4pt}
\begin{tabular}{lcccccc|c}
\toprule
& \multicolumn{3}{c}{Non-agent coding}
& \multicolumn{1}{c}{Agent OOD}
& \multicolumn{2}{c}{Tool use}
& \\
\cmidrule(lr){2-4} \cmidrule(lr){5-5} \cmidrule(lr){6-7}
Setting (Qwen2.5-Coder-14B-base)
& LiveCode & OJBench & FSB-EN
& Terminal
& $\tau$-bench & BFCL
& Avg \\
\midrule
Instruct (no agentic training)
  & 37.20 & 5.20 & 53.80
  & 0.00
  & 5.70 & 23.20
  & 20.85 \\
+ R2E-Gym
  & 24.10\,(\textcolor{red}{$\pm$?})
  & 2.80\,(\textcolor{red}{$\pm$?})
  & 47.72\,(\textcolor{red}{$\pm$?})
  & 1.28\,(\textcolor{red}{$\pm$?})
  & 0.40\,(\textcolor{red}{$\pm$?})
  & 15.70\,(\textcolor{red}{$\pm$?})
  & 15.33 \\
\textbf{+ FIM-Midtrain + R2E-Gym}
  & \textbf{35.20}\,(\textcolor{red}{$\pm$?})
  & \textbf{4.74}\,(\textcolor{red}{$\pm$?})
  & \textbf{48.07}\,(\textcolor{red}{$\pm$?})
  & \textbf{2.47}\,(\textcolor{red}{$\pm$?})
  & \textbf{4.40}\,(\textcolor{red}{$\pm$?})
  & \textbf{18.20}\,(\textcolor{red}{$\pm$?})
  & \textbf{18.85} \\
\rowcolor{blue!8}
$\Delta$ (vs.\ post-training only)
  & $+11.10$ & $+1.94$ & $+0.35$
  & $+1.19$
  & $+4.00$ & $+2.50$
  & $+3.52$ \\
\bottomrule
\end{tabular}
\end{table}

\paragraph{Agentic post-training has a substantial hidden capability cost.}
Applying R2E-Gym alone to the 14B base reduces every non-agent coding and
tool-use benchmark relative to the Instruct ceiling: LiveCodeBench drops
by $13.10$ points, BFCL by $7.50$, FullStackBench-EN by $6.08$,
$\tau$-bench by $5.30$, and OJBench by $2.40$. Averaged across the six
benchmarks the post-trained model loses $5.52$ points relative to the
Instruct ceiling. This regression is rarely highlighted in agent papers,
but it is the implicit cost being paid for SWE-Bench gains.

\paragraph{Mid-training largely closes the regression gap.}
Adding FIM mid-training before the same post-training restores
LiveCodeBench by $+11.10$, OJBench by $+1.94$ (within $0.46$ of the
Instruct ceiling), and FullStackBench-EN by $+0.35$. The six-benchmark
average rises from $15.33$ to $18.85$, an increase of $+3.52$ over the
post-training-only configuration, while the SWE-Bench-Verified gain on
the same base is preserved (Table~\ref{tab:main}). In other words,
mid-training improves the cost--benefit profile of agentic post-training:
the in-domain target metric improves, and the bulk of the off-distribution
capability erosion is undone in the same training run.

\paragraph{Cross-domain transfer to non-coding tool use.}
The most direct test of our motivating hypothesis lies in the rightmost
columns of Table~\ref{tab:cap}: $\tau$-bench and BFCL contain no Python
code-editing data, and our mid-training corpus contains no tool-use
trajectories. Mid-training nevertheless improves $\tau$-bench by
$+4.00$ and BFCL by $+2.50$ over post-training alone, with a directionally
consistent recovery on the also-OOD Terminal-Bench ($+1.19$). Because the
mid-training data carries no tool-use signal a priori, the only mechanism
by which it can affect these benchmarks is by installing a structural
prior at mid-training time that survives subsequent post-training. The
fact that this prior moves $\tau$-bench and BFCL in the same direction
as the in-domain SWE-Bench gains is consistent with our hypothesis that
the function-call site and the agent tool-call step share an underlying
conditioning structure (Section~\ref{sec:method:motivation}).


\subsection{Ablation Studies}
\label{sec:exp:ablation}

We conduct two controlled ablations, both on the Qwen2.5-Coder-7B base
with R2E-Gym post-training and a shared mid-training-free baseline. The
first isolates the role of the chain-of-thought rationale embedded in the
FIM middle span; the second isolates the role of the function-aware
selection pipeline. All selection variants in Block (B) operate under a
controlled budget: each variant selects exactly $100$K single functions
for masking, and basic hard filters (LoC $\in [10, 200]$, dunder methods
excluded) are applied uniformly. Differences across rows therefore reflect
\emph{which} functions are selected rather than \emph{how many}. Mask
granularity beyond single functions, score-threshold sensitivity, and
mid-training token budget are reserved for future investigation; due to
compute constraints, the ablations run at the 7B scale only.

% -----------------------------------------------------------------------------
% Ablation table
% -----------------------------------------------------------------------------
\begin{table}[t]
\centering
\caption{Ablation studies on the Qwen2.5-Coder-7B base with R2E-Gym
post-training. \textbf{(A)} varies the source of the rationale embedded in
the FIM middle span. \textbf{(B)} varies the function-selection algorithm
with all other components fixed (CoT source held to Gemini-3, budget
fixed to 100K single functions). The two blocks share the same baseline
(\emph{w/o mid-train}) and the same full configuration. Bold marks our
default recipe.}
\label{tab:ablation}
\small
\setlength{\tabcolsep}{8pt}
\begin{tabular}{lccc}
\toprule
Setting & SWE-Bench-Verified & SWE-Bench-Lite & Average \\
\midrule
\multicolumn{4}{l}{\emph{(A) Chain-of-thought source (selection fixed to ``full'')}} \\
\midrule
w/o mid-train (baseline)            & 15.00 & 11.33 & 13.17 \\
+ FIM, no CoT                       & 15.60 & 11.00 & 13.30 \\
+ FIM, self-CoT (Qwen2.5-Coder-7B)  & 16.40 & 13.30 & 14.85 \\
\textbf{+ FIM, Gemini-3 CoT (full)} & \textbf{17.80} & \textbf{15.00} & \textbf{16.40} \\
\midrule
\multicolumn{4}{l}{\emph{(B) Function-selection algorithm (CoT fixed to Gemini-3, budget fixed to 100K functions)}} \\
\midrule
w/o mid-train (baseline)               & 15.00 & 11.33 & 13.17 \\
Random                                 & \textcolor{red}{XX.XX} & \textcolor{red}{XX.XX} & \textcolor{red}{XX.XX} \\
Gemini-selected                        & \textcolor{red}{XX.XX} & \textcolor{red}{XX.XX} & \textcolor{red}{XX.XX} \\
PDG only                               & \textcolor{red}{XX.XX} & \textcolor{red}{XX.XX} & \textcolor{red}{XX.XX} \\
PDG + complexity ($\hat{H}$)           & \textcolor{red}{XX.XX} & \textcolor{red}{XX.XX} & \textcolor{red}{XX.XX} \\
PDG + inferability ($\hat{I}$)         & \textcolor{red}{XX.XX} & \textcolor{red}{XX.XX} & \textcolor{red}{XX.XX} \\
\textbf{Full (PDG + $\hat{H}$ + $\hat{I}$)} & \textbf{17.80} & \textbf{15.00} & \textbf{16.40} \\
\bottomrule
\end{tabular}
\end{table}

\paragraph{FIM structure contributes independently of CoT distillation.}
Block (A) of Table~\ref{tab:ablation} disarms a natural reviewer concern:
that our improvements stem primarily from distilling reasoning from a
frontier model. Removing the rationale entirely (\emph{no CoT}) still
improves the average by $+0.13$ over the no-mid-training baseline, with a
$+0.6$ gain on SWE-Bench-Verified, indicating that the function-aware FIM
objective by itself contributes some signal even without any external
rationale. Replacing Gemini-3 with rationales generated by the model
being trained (\emph{self-CoT}) reaches an average of $14.85$, recovering
$1.7$ of the $3.2$-point gap between \emph{no CoT} and \emph{Gemini-3 CoT};
high-quality external rationales help, but a sizable fraction of the
gain is reproducible without a frontier teacher.

\paragraph{The function-selection algorithm is the dominant lever within mid-training.}
Block (B) varies the selection algorithm with the CoT source held fixed
to Gemini-3 and the selection budget fixed to 100K single functions, so
that all rows incur identical mid-training cost.
\emph{Random} masking establishes the floor: with no structural priors
beyond the basic hard filters, mid-training reaches
\textcolor{red}{XX.XX} on average, a margin of \textcolor{red}{$+$X.XX}
over no mid-training. Letting Gemini-3 directly choose which functions
to mask (\emph{Gemini-selected})---a strong non-structural baseline that
uses the same teacher we use for rationales---reaches
\textcolor{red}{XX.XX}, indicating that frontier judgment on \emph{which}
function to mask is helpful but not by itself sufficient. Restricting the
candidate pool to functions with at least one neighbor in the program
dependency graph (\emph{PDG only}) reaches \textcolor{red}{XX.XX}, a
further \textcolor{red}{$+$X.XX} over Random: functions embedded in
caller-callee or sibling structure are intrinsically more useful as FIM
targets than isolated ones, even before any scoring is applied. Adding a
scoring function on top of the PDG pool yields the next layer of gain:
\emph{PDG + complexity} reaches \textcolor{red}{XX.XX} and
\emph{PDG + inferability} reaches \textcolor{red}{XX.XX}, showing that
both the intrinsic difficulty of the target and its recoverability from
surrounding context independently contribute beyond the structural
filter. Combining the two scores via the harmonic-mean-like
$\hat{H}\,\hat{I}/(\hat{H}+\hat{I})$ form---the \emph{full}
configuration---reaches $16.40$, a further \textcolor{red}{$+$X.XX} over
the better single-score variant. The combination is therefore not
redundant: $\hat{H}$ and $\hat{I}$ measure complementary properties of a
candidate, and selecting targets that score highly on only one of the
two yields either unlearnable noise (high complexity, low inferability)
or trivial fills (low complexity, high inferability).